\subsubsection{Ergodicité}

Une chaîne de Markov est dite ergodique si elle est irréductible 
(il est possible d'atteindre n'importe quel état depuis n'importe quel autre) et apériodique. 
Pour une telle chaîne, le théorème ergodique stipule que la proportion du temps passé dans un 
état $i$ sur une trajectoire infiniment longue converge vers la probabilité 
stationnaire $\pi_i$ de cet état :
\begin{equation}
    \lim_{N \to \infty} \frac{1}{N} \sum_{t=1}^N \mathbf{1}_{\{X_t = i\}} = \pi_i
\end{equation}

Pour vérifier expérimentalement cette propriété, nous avons simulé une trajectoire de la 
chaîne en générant une nouvelle séquence artificielle de $N = 50\,000$ nucléotides. À 
chaque étape, la transition vers le nucléotide suivant a été tirée aléatoirement en respectant 
strictement les probabilités de la matrice de transition $P$ estimée à la Question 1.1.1. 
Nous avons ensuite calculé les fréquences empiriques d'apparition de chaque nucléotide dans 
cette séquence générée.

Le tableau ci-dessous compare ces fréquences empiriques temporelles avec la distribution 
stationnaire $\pi_\infty$ calculée analytiquement à la Question 1.1.3 :

\begin{center}
\renewcommand{\arraystretch}{1.2} % Donne un peu d'espace dans les lignes du tableau
\begin{tabular}{|c|c|c|}
    \hline
    \textbf{Nucléotide} & \textbf{Fréquence empirique (N=50 000)} & \textbf{Distribution stationnaire ($\pi_\infty$)} \\
    \hline
    A & 0.2879 & 0.2881 \\ 
    C & 0.3483 & 0.3497 \\ 
    G & 0.2117 & 0.2091 \\ 
    T & 0.1521 & 0.1531 \\ 
    \hline
\end{tabular}
\end{center}

\subparagraph{Conclusion}
Nous constatons que les fréquences empiriques observées sur une longue réalisation de la 
chaîne sont extrêmement proches des probabilités de la distribution stationnaire $\pi_\infty$ 
(aux fluctuations statistiques d'échantillonnage près). Cela confirme numériquement que 
notre modèle de séquence d'ADN forme bien une chaîne de Markov ergodique : 
l'exploration temporelle du système (simulée sur 50 000 étapes) reflète fidèlement son 
équilibre spatial théorique.